\documentclass[a4paper,11pt]{article}
\usepackage[utf8]{inputenc}
\usepackage{geometry}
\usepackage{xcolor}
\usepackage{listings}
\usepackage{hyperref}
\usepackage{graphicx}
\usepackage{amsmath}
\usepackage{amssymb}
\usepackage{tcolorbox}
\usepackage{float}
\usepackage{booktabs}
\usepackage{tikz}
\usetikzlibrary{arrows.meta, positioning, shapes.geometric, matrix}

% ── Page geometry ──────────────────────────────────────────────────────────────
\geometry{top=2.5cm, bottom=2.5cm, left=2.5cm, right=2.5cm}

% ── Custom colours ─────────────────────────────────────────────────────────────
\definecolor{ibmblue}{RGB}{0,45,156}
\definecolor{ibmred}{RGB}{250,75,76}
\definecolor{codegreen}{rgb}{0,0.6,0}
\definecolor{codegray}{rgb}{0.5,0.5,0.5}
\definecolor{codepurple}{rgb}{0.58,0,0.82}
\definecolor{backcolour}{rgb}{0.95,0.95,0.92}
\definecolor{tealcolor}{RGB}{0,128,128}
\definecolor{orangecolor}{RGB}{230,100,0}
\definecolor{goldcolor}{RGB}{180,140,0}

% ── Colour helpers ─────────────────────────────────────────────────────────────
\newcommand{\crd}[1]{\textcolor{ibmred}{#1}}
\newcommand{\cpur}[1]{\textcolor{codepurple}{#1}}
\newcommand{\cblu}[1]{\textcolor{ibmblue}{#1}}
\newcommand{\cgr}[1]{\textcolor{codegreen}{#1}}
\newcommand{\corg}[1]{\textcolor{orangecolor}{#1}}
\newcommand{\cteal}[1]{\textcolor{tealcolor}{#1}}
\newcommand{\cgold}[1]{\textcolor{goldcolor}{#1}}

% ── Listings ───────────────────────────────────────────────────────────────────
\lstdefinestyle{mystyle}{
    backgroundcolor=\color{backcolour},
    commentstyle=\color{codegreen},
    keywordstyle=\color{magenta},
    numberstyle=\tiny\color{codegray},
    stringstyle=\color{codepurple},
    basicstyle=\ttfamily\footnotesize,
    breakatwhitespace=false, breaklines=true,
    captionpos=b, keepspaces=true,
    numbers=left, numbersep=5pt,
    showspaces=false, showstringspaces=false,
    showtabs=false, tabsize=2
}
\lstset{style=mystyle}

% ── tcolorbox environments ─────────────────────────────────────────────────────
\newtcolorbox{studentbox}{colback=blue!5!white,colframe=blue!75!black,
    title=Student Activity, fonttitle=\bfseries}
\newtcolorbox{warningbox}[1][Key Concept]{%
    colback=yellow!10!white,colframe=orange!80!black,
    title=#1, fonttitle=\bfseries}
\newtcolorbox{mathbox}[1][Derivation]{%
    colback=red!5!white,colframe=red!75!black,
    title=#1, fonttitle=\bfseries}
\newtcolorbox{defbox}[1][Definition]{%
    colback=teal!5!white,colframe=teal!70!black,
    title=#1, fonttitle=\bfseries}
\newtcolorbox{finbox}[1][Financial Interpretation]{%
    colback=green!5!white,colframe=green!65!black,
    title=#1, fonttitle=\bfseries}
\newtcolorbox{bridgebox}{colback=purple!5!white,colframe=purple!70!black,
    title=Connection to Lab 4, fonttitle=\bfseries}

% ── Math shorthands ────────────────────────────────────────────────────────────
\newcommand{\ket}[1]{|#1\rangle}
\newcommand{\bra}[1]{\langle#1|}
\newcommand{\expect}[1]{\langle #1 \rangle}
\newcommand{\Mod}[1]{\ (\mathrm{mod}\ #1)}
\newcommand{\tr}{\mathrm{Tr}}

\title{\textbf{Laboratory Guide 5: Quantum Portfolio Optimization}\\
\large Finance \& Industry Applications of Variational Quantum Algorithms}
\author{Course: Quantum Computing Foundation}
\date{}

\begin{document}

\maketitle
\tableofcontents
\newpage

% ══════════════════════════════════════════════════════════════════════════════
\section{Introduction and Objectives}
% ══════════════════════════════════════════════════════════════════════════════

In Lab 4 you built the full technical machinery of variational quantum
algorithms: the parameter-shift rule, VQE, and QAOA applied to the MaxCut
combinatorial problem. This laboratory is the direct application of that
machinery to a real-world industrial problem: \textbf{portfolio optimization
in quantitative finance}.

The core mathematical structure is identical to Lab 4 — a Quadratic
Unconstrained Binary Optimization (QUBO) problem mapped to a QAOA Hamiltonian.
The novelty is threefold: (i) the QUBO is derived from real financial data,
(ii) interpreting the output requires both a technical and a business lens,
and (iii) the problem has a natural classical solution that allows direct
benchmarking.

\vspace{0.3cm}
\noindent\textbf{Prerequisites:}
\begin{itemize}
    \item Lab 4: Variational Quantum Algorithms, QAOA, QUBO formulation
    \item Basic probability and statistics: mean, variance, covariance
    \item No prior knowledge of quantitative finance is assumed
\end{itemize}

\noindent\textbf{Learning Objectives:}
\begin{enumerate}
    \item Model a combinatorial portfolio selection problem as a QUBO.
    \item Map the QUBO to a quantum Hamiltonian and build the QAOA circuit.
    \item Run QAOA on a 4-asset universe (AAPL, MSFT, JPM, XOM) and
          verify the result against brute-force.
    \item Interpret results from both a business and a technical perspective.
    \item Understand the scaling advantage for large $n$ and the challenges
          of NISQ hardware.
\end{enumerate}

\begin{warningbox}[Why Quantum for Portfolio Optimization?]
For $n$ assets selecting $K$, the number of candidate portfolios is
$\binom{n}{K}$. Classical branch-and-bound solvers handle this well up
to $n \approx 1000$. However, adding realistic constraints (transaction
costs, sector limits, ESG scores, integer lots) makes the problem
NP-hard in general. QAOA provides a \emph{near-term heuristic} that
may achieve quantum advantage for $n \gtrsim 100$ constrained assets —
a scale that matters for institutional portfolios (e.g.\ Russell 2000
universe, $n=2000$).
\end{warningbox}

% ══════════════════════════════════════════════════════════════════════════════
\section{Part A: Classical Portfolio Theory}
% ══════════════════════════════════════════════════════════════════════════════

\subsection{Mean-Variance Optimization (Markowitz, 1952)}

Let $x \in \mathbb{R}^n$ be a portfolio weight vector with
$\sum_i x_i = 1$ (continuous). The mean-variance problem is:

\begin{equation}
    \min_x\; \lambda\, x^T \Sigma\, x - \boldsymbol\mu^T x,
    \label{eq:mv-cont}
\end{equation}

where $\Sigma \in \mathbb{R}^{n\times n}$ is the asset return
\emph{covariance matrix}, $\boldsymbol\mu \in \mathbb{R}^n$ is the
vector of \emph{expected returns}, and $\lambda > 0$ is the
\emph{risk-aversion} parameter. The first term penalises risk (portfolio
variance), the second rewards expected return.

\begin{defbox}[Risk-Aversion Parameter $\lambda$]
\begin{itemize}
    \item $\lambda \to 0$: maximise returns regardless of risk
          (aggressive investor).
    \item $\lambda \to \infty$: minimise variance regardless of returns
          (defensive investor, cash-like).
    \item Typical institutional values: $\lambda \in [0.5,\, 5]$.
\end{itemize}
For this lab we use $\lambda = 1$ as the baseline.
\end{defbox}

\subsection{The Combinatorial Version and its Hardness}

In practice, portfolio managers must select \emph{exactly $K$ assets}
from a universe of $n$ candidates and invest equally. The binary version
of equation~\eqref{eq:mv-cont} is:

\begin{equation}
    \min_{x \in \{0,1\}^n}\; \lambda\, x^T \Sigma\, x
    - \boldsymbol\mu^T x,
    \quad\text{subject to}\quad \mathbf{1}^T x = K.
    \label{eq:mv-binary}
\end{equation}

For $n=4$, $K=2$ there are only $\binom{4}{2} = 6$ candidates — trivially
enumerable. For $n=50$, $K=10$ there are $\binom{50}{10} \approx 10^{10}$
candidates. For $n=500$, $K=50$: $\sim\!10^{82}$ — far beyond brute-force.
The continuous relaxation (equation~\eqref{eq:mv-cont}) gives the
\emph{efficient frontier} but loses the discreteness that real portfolios
require.

\subsection{Asset Universe: AAPL, MSFT, JPM, XOM (2020--2024)}

We work with the four assets in Table~\ref{tab:assets}. The covariance
matrix $\Sigma$ and expected returns $\boldsymbol\mu$ are computed from
daily closing prices over 2020--2024 (annualised).

\begin{table}[H]
\centering
\caption{Asset data used in this laboratory (annualised, 2020--2024).}
\label{tab:assets}
\begin{tabular}{clccc}
    \toprule
    Index & Ticker & Sector & $\mu_i$ & $\sigma_i = \sqrt{\Sigma_{ii}}$ \\
    \midrule
    0 & AAPL  & Technology    & 0.340 & 0.290 \\
    1 & MSFT  & Technology    & 0.270 & 0.270 \\
    2 & JPM   & Financials    & 0.150 & 0.240 \\
    3 & XOM   & Energy        & 0.220 & 0.350 \\
    \bottomrule
\end{tabular}
\end{table}

\noindent The annualised covariance matrix (approximate):
\begin{equation}
    \Sigma = \begin{pmatrix}
        \cgr{0.0841} & 0.0588 & 0.0313 & 0.0203 \\
        0.0588 & \cgr{0.0729} & 0.0272 & 0.0170 \\
        0.0313 & 0.0272 & \cgr{0.0576} & 0.0294 \\
        0.0203 & 0.0170 & 0.0294 & \cgr{0.1225}
    \end{pmatrix}.
    \label{eq:sigma}
\end{equation}
The diagonal entries $\Sigma_{ii} = \sigma_i^2$ are the individual asset
variances; the off-diagonal entries capture return correlations.
AAPL--MSFT have the highest co-movement (both tech-sector), while
XOM (energy) is weakly correlated with the others.

% ══════════════════════════════════════════════════════════════════════════════
\section{Part B: QUBO Formulation}
% ══════════════════════════════════════════════════════════════════════════════

\subsection{From Constrained to Unconstrained}

We convert the equality constraint $\sum_i x_i = K$ into a quadratic
penalty added to the objective:

\begin{equation}
    \min_{x\in\{0,1\}^n}\;
    \underbrace{\lambda\, x^T\Sigma x - \boldsymbol\mu^T x}_{\text{cost}}
    \;+\;
    \underbrace{\lambda_\text{pen}\!\left(\sum_{i=1}^n x_i - K\right)^{\!2}}_{\text{budget penalty}},
    \label{eq:qubo-full}
\end{equation}

where $\lambda_\text{pen}$ must be large enough that any infeasible
solution (wrong number of assets) is always worse than any feasible one.
A safe choice is $\lambda_\text{pen} > \max_{x,\text{feasible}} |C(x)|$.

Expanding and collecting, equation~\eqref{eq:qubo-full} takes the
standard QUBO form:

\begin{equation}
    \min_{x\in\{0,1\}^n}\; x^T Q\, x,
    \label{eq:qubo}
\end{equation}

\noindent with:
\begin{equation}
    Q_{ij} = \begin{cases}
        \lambda\,\Sigma_{ii} - \mu_i + \lambda_\text{pen}(1-2K)
        & i = j \\[4pt]
        \lambda\,\Sigma_{ij} + 2\lambda_\text{pen}
        & i \neq j.
    \end{cases}
    \label{eq:Q}
\end{equation}

\begin{mathbox}[QUBO Matrix for $n=4$, $K=2$, $\lambda=1$, $\lambda_\text{pen}=3$]
Substituting equation~\eqref{eq:sigma} and $\mu = (0.34, 0.27, 0.15, 0.22)$
into equation~\eqref{eq:Q} with $\lambda_\text{pen}(1-2K) = 3(1-4) = -9$:
\[
    Q = \lambda\Sigma - \mathrm{diag}(\boldsymbol\mu)
    + \lambda_\text{pen}
    \begin{pmatrix}
        1-2K & 2 & 2 & 2 \\
        2 & 1-2K & 2 & 2 \\
        2 & 2 & 1-2K & 2 \\
        2 & 2 & 2 & 1-2K
    \end{pmatrix}
\]
\[
    Q = \begin{pmatrix}
        -\crd{8.8259} & 6.1764 & 6.0939 & 6.0609 \\
        6.1764 & -\crd{8.9571} & 6.0816 & 6.0510 \\
        6.0939 & 6.0816 & -\crd{9.0324} & 6.0882 \\
        6.0609 & 6.0510 & 6.0882 & -\crd{8.7575}
    \end{pmatrix}.
\]
The large off-diagonal terms reflect the strong budget penalty;
the negative diagonal terms reflect the return incentive.
\end{mathbox}

\subsection{Classical Brute-Force Baseline}

For $n=4$, $K=2$ we can enumerate all $\binom{4}{2}=6$ portfolios
exactly. Table~\ref{tab:bf} shows the cost $x^T Qx$ and the
corresponding financial interpretation.

\begin{table}[H]
\centering
\caption{Brute-force enumeration of all 2-asset portfolios ($\lambda=1$).}
\label{tab:bf}
\begin{tabular}{clcc}
    \toprule
    Bitstring $x$ & Assets selected & Cost $x^TQx$ & Rank \\
    \midrule
    $\cblu{1100}$ & \cblu{AAPL + MSFT} & $\cblu{-0.335}$ & \cblu{1 (optimal)} \\
    $1001$ & AAPL + XOM        & $-0.313$ & 2 \\
    $1010$ & AAPL + JPM        & $-0.286$ & 3 \\
    $0101$ & MSFT + XOM        & $-0.261$ & 4 \\
    $0110$ & MSFT + JPM        & $-0.235$ & 5 \\
    $0011$ & JPM  + XOM        & $-0.131$ & 6 \\
    \bottomrule
\end{tabular}
\end{table}

\begin{finbox}
\textbf{Why AAPL + MSFT?} Despite having the highest inter-asset
covariance (both technology sector), AAPL and MSFT also have the
highest expected returns ($0.34$ and $0.27$ respectively). For
$\lambda=1$ the return advantage outweighs the diversification cost.
As $\lambda$ increases (more risk-averse investor), the optimal
portfolio shifts toward AAPL + XOM (energy provides diversification)
and eventually toward JPM + XOM (defensive, low correlation).
\end{finbox}

% ══════════════════════════════════════════════════════════════════════════════
\section{Part C: Quantum Hamiltonian and QAOA Circuit}
% ══════════════════════════════════════════════════════════════════════════════

\subsection{QUBO $\to$ Ising Hamiltonian}

Substituting $x_i = (1 - Z_i)/2$ into equation~\eqref{eq:qubo}:

\begin{align}
    x^T Q\, x &= \sum_{i,j} Q_{ij} x_i x_j
    = \sum_{i,j} Q_{ij} \frac{(1-Z_i)(1-Z_j)}{4} \notag\\
    &= \frac{1}{4}\sum_{i,j} Q_{ij}
    \bigl(I - Z_i - Z_j + Z_iZ_j\bigr) \notag\\
    &\equiv H_\text{portfolio} \;+\; \text{const.}
    \label{eq:ising-map}
\end{align}

Collecting terms, the cost Hamiltonian is:

\begin{equation}
    H_\text{portfolio}
    = \sum_{i<j} \frac{Q_{ij}}{2}\, Z_iZ_j
    + \sum_i \left[-\frac{1}{2}\sum_j Q_{ij}\right] Z_i
    + \text{const.}
    \label{eq:H-portfolio}
\end{equation}

\begin{bridgebox}
The structure of $H_\text{portfolio}$ is a sum of $Z_iZ_j$ and $Z_i$
Pauli operators — \emph{exactly} the same form as the MaxCut Hamiltonian
in Lab 4. The QAOA circuit implementation is therefore identical:
\begin{itemize}
    \item $e^{-i\gamma Z_i Z_j}$: CNOT\,--\,$R_z(2\gamma)$\,--\,CNOT
    \item $e^{-i\gamma Z_i}$: $R_z(2\gamma)$ directly
    \item $e^{-i\beta H_B}$: $R_x(2\beta)$ on all qubits
\end{itemize}
Only the coefficients in front of the Pauli strings change.
\end{bridgebox}

\subsection{Explicit Hamiltonian Terms ($n=4$, $\lambda=1$)}

Applying equation~\eqref{eq:H-portfolio} to our $Q$ matrix yields
the following non-trivial Pauli terms (omitting the constant):

\begin{center}
\begin{tabular}{lrl}
    \toprule
    Pauli term & Coefficient & Origin \\
    \midrule
    $Z_0 Z_1$ & $+3.0882$ & AAPL--MSFT covariance + penalty \\
    $Z_0 Z_2$ & $+3.0470$ & AAPL--JPM covariance + penalty \\
    $Z_0 Z_3$ & $+3.0305$ & AAPL--XOM covariance + penalty \\
    $Z_1 Z_2$ & $+3.0408$ & MSFT--JPM covariance + penalty \\
    $Z_1 Z_3$ & $+3.0255$ & MSFT--XOM covariance + penalty \\
    $Z_2 Z_3$ & $+3.0441$ & JPM--XOM covariance + penalty \\
    $Z_0$     & $+0.1568$ & AAPL diagonal term \\
    $Z_1$     & $+0.0929$ & MSFT diagonal term \\
    $Z_2$     & $-0.0474$ & JPM diagonal term \\
    $Z_3$     & $+0.1238$ & XOM diagonal term \\
    \bottomrule
\end{tabular}
\end{center}

\subsection{QAOA Circuit Structure ($p=1$)}

The $p=1$ QAOA circuit alternates one cost unitary and one mixer:

\begin{equation}
    \ket{\gamma,\beta} =
    e^{-i\beta\sum_i X_i}\,
    e^{-i\gamma H_\text{portfolio}}\,
    H^{\otimes 4}\ket{0000}.
\end{equation}

The cost unitary decomposes into one $R_z$ gate per $Z_i$ term and
one CNOT\,--\,$R_z$\,--\,CNOT block per $Z_iZ_j$ term. For $n=4$:
4 single-qubit + $4 \times 2 = 8$ two-qubit blocks (for $\binom{4}{2}=6$ edges)
= \textbf{4 $R_z$ + 6 CNOT pairs} per cost layer. Total circuit depth
for $p=1$: $\mathcal{O}(n^2)$ gates.

% ══════════════════════════════════════════════════════════════════════════════
\section{Part D: Result Interpretation}
% ══════════════════════════════════════════════════════════════════════════════

\subsection{Reading the Output Distribution}

After QAOA optimisation, we sample the circuit $M$ times and obtain a
histogram over $2^n = 16$ bitstrings. The \emph{technical} interpretation:
the state $\ket{1100}$ (AAPL=1, MSFT=1, JPM=0, XOM=0, in Qiskit
little-endian) should have the highest probability.

\begin{finbox}[Technical $\to$ Business Translation]
\begin{center}
\begin{tabular}{cl}
    \toprule
    Bitstring & Business meaning \\
    \midrule
    $\ket{1100}$ & Invest in AAPL and MSFT \\
    $\ket{1001}$ & Invest in AAPL and XOM \\
    $\ket{0011}$ & Invest in JPM and XOM \\
    $\ket{0000}$ & Infeasible (no assets) \\
    $\ket{1111}$ & Infeasible (all assets, $K \neq 2$) \\
    \bottomrule
\end{tabular}
\end{center}
Infeasible states have high energy due to the penalty term —
they should appear with near-zero probability after optimisation.
\end{finbox}

\subsection{Approximation Quality and Business Risk}

QAOA with $p=1$ does not always find the exact optimum. Define the
\emph{success probability} $P^* = P(\ket{1100})$: the probability
the circuit returns the exact optimal portfolio. For $n=4$, $p=1$:
$P^* \approx 40$--$60\%$. For production use, we would:

\begin{enumerate}
    \item Sample the circuit $M$ times (e.g.\ $M=1000$).
    \item Take the most frequent \emph{feasible} bitstring as the answer.
    \item Compare its cost against the top-$k$ classical candidates.
\end{enumerate}

Increasing $p$ improves $P^*$ but also increases circuit depth and
susceptibility to hardware noise.

\subsection{Sensitivity Analysis: Risk Aversion $\lambda$}

\begin{mathbox}[How $\lambda$ Shifts the Optimal Portfolio]
At $\lambda = 0.5$ (return-focused): $\cblu{\text{AAPL}+\text{MSFT}}$ optimal. \\
At $\lambda = 1.0$ (balanced):      $\cblu{\text{AAPL}+\text{MSFT}}$ optimal. \\
At $\lambda = 2.0$ (risk-averse):   $\corg{\text{AAPL}+\text{XOM}}$ may emerge. \\
At $\lambda = 5.0$ (very defensive): $\crd{\text{JPM}+\text{XOM}}$ (lowest covariance
                                                                    pair) wins.

\noindent The quantum circuit must be re-optimised for each $\lambda$
because the Hamiltonian $H_\text{portfolio}(\lambda)$ changes.
The optimal QAOA parameters $(\gamma^*,\beta^*)$ also shift.
\end{mathbox}

% ══════════════════════════════════════════════════════════════════════════════
\section{Part E: Qiskit Python Implementation}
% ══════════════════════════════════════════════════════════════════════════════

The full implementation is in the companion notebook. Three exercises
are provided:

\begin{enumerate}
    \item \textbf{QUBO Construction:} Build $Q$ from real asset data,
          enumerate all $\binom{4}{2}$ portfolios classically, identify
          the optimum.
    \item \textbf{QAOA Optimisation:} Run QAOA with $p \in \{1,2\}$,
          compare the measured distribution against brute-force.
    \item \textbf{Business Analysis:} Plot the efficient frontier,
          extract the top-3 most probable portfolios from the QAOA
          sample, and compute their actual risk-return profiles.
\end{enumerate}

\begin{tcolorbox}[colback=green!5!white,colframe=green!75!black,
    title=Key Code Structure, fonttitle=\bfseries]
\begin{lstlisting}[language=Python]
# Core objects: same structure as Lab 4, different Hamiltonian
def build_portfolio_hamiltonian(Q: np.ndarray) -> SparsePauliOp:
    n     = Q.shape[0]
    terms = []
    for i in range(n):
        for j in range(i+1, n):
            coeff = Q[i,j] / 2           # ZiZj coefficient
            terms.append((zz_str(i,j,n), coeff))
        z_coeff = -sum(Q[i,:]) / 2       # Zi coefficient
        terms.append((z_str(i,n), z_coeff))
    return SparsePauliOp.from_list(terms).simplify()

# QAOA circuit: identical to Lab 4's build_qaoa_circuit
# Only H_C = H_portfolio changes; circuit structure is the same
\end{lstlisting}
\end{tcolorbox}

% ══════════════════════════════════════════════════════════════════════════════
\section{Part F: Scaling and Outlook}
% ══════════════════════════════════════════════════════════════════════════════

\subsection{Scaling to Larger Portfolios}

\begin{table}[H]
\centering
\caption{QAOA circuit scaling for portfolio optimization.}
\label{tab:scaling}
\begin{tabular}{cccccc}
    \toprule
    $n$ & $\binom{n}{K}$ & QAOA qubits & ZZ terms & Classical gates (p=1) \\
    \midrule
    4  & 6          & 4   & 6   & 16  \\
    6  & 15         & 6   & 15  & 36  \\
    10 & 252        & 10  & 45  & 100 \\
    20 & 184,756    & 20  & 190 & 400 \\
    50 & $\sim 10^{12}$ & 50  & 1,225 & 2,500 \\
    \bottomrule
\end{tabular}
\end{table}

The QAOA circuit grows as $\mathcal{O}(n^2)$ gates (from the $\binom{n}{2}$
ZZ interaction terms), while the classical search space grows as
$\binom{n}{K} \sim n^K / K!$. For $n=20$, $K=5$: classical brute-force
checks $\sim\!10^5$ portfolios (fast classically); for $n=50$, $K=10$:
$\sim\!10^{12}$ (intractable). The QAOA circuit for $n=50$ requires
only $1225$ two-qubit gates per layer — within reach of near-term hardware.

\subsection{Limitations and Open Questions}

\begin{itemize}
    \item \textbf{Barren plateaus:} For $n \gtrsim 50$, the QAOA landscape
          may be exponentially flat, requiring problem-informed initial
          parameters.
    \item \textbf{Hardware noise:} The $\sim\!1225$ CX gates at $n=50$
          with $\sim\!0.5\%$ error each give an overall fidelity of
          $0.995^{1225} \approx 0.002$ — noise dominates.
    \item \textbf{Quantum advantage:} Not yet demonstrated for portfolio
          optimization. Active research area (JPMorgan, Goldman Sachs,
          IBM Research).
    \item \textbf{Beyond QUBO:} Transaction costs, lot sizes, turnover
          constraints are not naturally quadratic — require slack variables
          or penalty engineering.
\end{itemize}

% ══════════════════════════════════════════════════════════════════════════════
\section{Part G: Assessment}
% ══════════════════════════════════════════════════════════════════════════════

\textbf{Q1.\ QUBO Formulation}\\
Why is the budget constraint $\sum_i x_i = K$ added as a penalty term
$\lambda_\text{pen}(\sum_i x_i - K)^2$ rather than enforced directly in
the quantum circuit?
\begin{itemize}
    \item[A)] Because quantum circuits cannot measure the number of 1s
              in a bitstring.
    \item[B)] Because the QAOA Hamiltonian must be a sum of Pauli operators
              (diagonal in the $Z$-basis), and a soft penalty converts the
              hard constraint into such a form.
    \item[C)] Because the penalty term is always zero for feasible solutions,
              making it costless.
\end{itemize}

\vspace{0.3cm}
\textbf{Q2.\ Hamiltonian Terms}\\
In the portfolio Hamiltonian $H_\text{portfolio}$, what is the physical
origin of the $Z_iZ_j$ terms?
\begin{itemize}
    \item[A)] The expected returns $\mu_i$ of the individual assets.
    \item[B)] The covariance $\Sigma_{ij}$ between assets $i$ and $j$,
              plus the budget penalty terms that couple all pairs of qubits.
    \item[C)] The CNOT gates in the QAOA mixer.
\end{itemize}

\vspace{0.3cm}
\textbf{Q3.\ Result Interpretation}\\
A QAOA circuit for the 4-asset portfolio returns the histogram:
$P(\ket{1100})=0.45$, $P(\ket{1001})=0.22$, $P(\ket{0011})=0.15$,
all others $< 0.05$. What is the recommended decision, and why is the
second-most-probable state a useful output?
\begin{itemize}
    \item[A)] Always select $\ket{1100}$; all other states are errors.
    \item[B)] Select $\ket{1100}$ as primary; $\ket{1001}$ is a useful
              second-best feasible portfolio that can serve as a
              diversification alternative.
    \item[C)] Re-run the circuit until one state reaches $>99\%$ probability.
\end{itemize}

\vspace{0.3cm}
\textbf{Q4.\ Scaling}\\
For $n=50$ assets and $K=10$, the classical brute-force requires
$\binom{50}{10} \approx 10^{10}$ cost evaluations. How many two-qubit
gates does the QAOA circuit for $n=50$ require per layer $p$?
\begin{itemize}
    \item[A)] $n = 50$
    \item[B)] $n^2 = 2500$
    \item[C)] $\binom{n}{2} = 1225$ (one CNOT pair per ZZ interaction term)
\end{itemize}

\vspace{1cm}
\hrule
\vspace{0.2cm}
\footnotesize{\textbf{Solutions:}
Q1: B (the Ising Hamiltonian is $\sum_i h_i Z_i + \sum_{ij} J_{ij} Z_iZ_j$;
    the penalty converts the constraint into this Pauli-operator form);
Q2: B (the $Q_{ij}$ for $i\neq j$ contains $\lambda\Sigma_{ij}$ from
    the risk term plus $2\lambda_\text{pen}$ from the penalty, both
    contributing to the $Z_iZ_j$ coefficient after the substitution
    $x_i=(1-Z_i)/2$);
Q3: B (in real applications the QAOA output is a probability distribution
    over solutions; the second-most-probable feasible state is a legitimate
    portfolio recommendation, not noise — it represents the second-best
    risk-return trade-off found by the algorithm);
Q4: C ($\binom{50}{2} = 1225$ ZZ terms each requiring one CNOT pair,
    giving $2\times 1225 = 2450$ CNOT gates per layer, plus 50 $R_z$
    gates; total $\mathcal{O}(n^2)$).}

\end{document}
